\scanpage{16}
Propriétés $(\unsure{P|\mathscr P}_1)$ et $(P_2)$ \unsure{pour} des faisceaux d'ensembles.

\deleted{\unsure{Proposition}}

\uline{\unsure{Définition} 1} \unsure{Soient|Soit} $X$ espace topologique, $Y$ une partie fermée
de $X$, \illegible{2} faisceau d'ensembles sur $X$. On dit
que $F$ \unsure{a} la propriété $(P_1)$ (resp. $(P_2)$)
\deleted{\unsure{rel. à}}\inserted{rel. à} \deleted{\unsure{un}} $Y$, si l'\unsure{hom} \unsure{canonique}
\[
F\longrightarrow i_*i^*(F)=i_*(F\mid U)
\]
(où $i:U=X-Y\longrightarrow X$ \unsure{est} l'\unsure{imm.} canonique)
est injective (resp. bijective).
\margin{N.B. \unsure{regarder|signaler} aussi la propriété: \unsure{surjective|surjection}....}
% slanted

(Cela signifie \unsure{donc} aussi que
pour tout ouvert $V\subset X$, l'application \deleted{\unsure{de}} restriction
\[
\Gamma(V,F)\unsure{-|\longrightarrow}\Gamma(V\cap U,F)
\]
est injective (resp. bijective).

\unsure{Appelons} \uline{espace noethérien local} un espace
noethérien dont \unsure{tte} partie fermée \unsure{irred} a un
pt générique et un seul\inserted{ [``espace noethérien \unsure{spectral|spécial}'']}, et tel que $X$ ait un
seul pt fermé. Définition de $\deleted{\unsure{X_x}}\operatorname{Loc}_x(X)$, $x$ pt d'un
espace loc. noethérien \unsure{spectral}. Définition de $\deleted{\unsure{F_x}}\operatorname{Loc}_x(F)$

\uline{Définition} 2 Soit $F$ un faisceau sur un espace
noethérien local $X$. On dit que $F$ \unsure{a} la
propriété $(P_1)$, resp. $(P_2)$, si \deleted{\unsure{le}} $F$ a la propriété
$(P_1)$ resp. $(P_2)$ rel. à $Y=\{x\}$, où $x$ est le pt fermé.
\margin{Cela signifie au \unsure{moins}, si $\deleted{\unsure{X}}U=X-\{x\}$, que $\Gamma(X,F)=F_x\longrightarrow\Gamma(U,\unsure{\Gamma|F})$ est injectif (resp. bijectif)}
% margin note written sideways
% Long vertical guides mark the two definitions. A bracket joins the sideways note to the foot of the page.
% archive sheet number: 15

